\documentclass{beamer}
\usefonttheme[onlymath]{serif}
% does not look nice, try deleting the line with the fontenc.
\usepackage[english]{babel}
\usepackage{amsmath}
\usepackage[latin1]{inputenc}
\usepackage{units}
\usepackage{colortbl}
\usepackage{multimedia}
\usepackage{bm}

\mode<presentation>
{
  \usetheme{Boadilla}
  \useoutertheme{infolines}
  \setbeamercovered{transparent} 
}
	

\title[Efficient Calibration of Embedded MPC]{\textsc{Efficient Calibration of Embedded MPC}}

%\subtitle{Industrial-scale experimental results\\} % (optional)

\author[]{Marco Forgione\inst{1}, Dario Piga\inst{1}, 
and Alberto Bemporad \inst{2}}

\institute[IDSIA]{
	\inst{1}IDSIA Dalle Molle Institute for Artificial Intelligence SUPSI-USI, Lugano, Switzerland \and
	\inst{2}IMT School for Advanced Studies Lucca, Lucca, Italy
	} 


\date[IFAC 2020]{IFAC 2020 World Congress\\ Berlin, Germany}


\subject{Performance-oriented MPC tuning}


%% MATH DEFINITIONS %%
%\newcommand{\R}{\mathcal{R}}
%\newcommand{\du}{\delta u}
%\newcommand{\dy}{\delta y}
%\newcommand{\DU}{\Delta U}
%\newcommand{\DY}{\Delta Y}
%\newcommand{\abs}[1]{\left | #1 \right |}
\newcommand{\norm}[1]{\left \lVert #1 \right \rVert}
%\newcommand{\relphantom}[1]{\mathrel{\phantom{#1}}}
%\newenvironment{matrixc}{\begin{array}{c}\left[}{\end{array}\right]}
\DeclareMathOperator*\argmin{arg \, min}
%\DeclareMathOperator*\argmax{arg \, max}
%\DeclareMathOperator*\fit{fit}
%\DeclareMathOperator*\RMSE{RMSE}
%\DeclareMathOperator*\diag{diag}
%\DeclareMathOperator*\diet{diet}
%\DeclareMathOperator*\Risk{Risk}
%\DeclareMathOperator*\Num{Num}
%\DeclareMathOperator*\Den{Den}
%\DeclareMathOperator*\Rat{Rat}
\DeclareMathOperator*\cov{cov}
%\DeclareMathOperator*\Var{Var}
%\DeclareMathOperator*\SSR{SSR}
%\setcounter{MaxMatrixCols}{20}
%\newcommand{\pdiff}[2]{\frac{\partial #1}{\partial #2}}
%\definecolor{mypink1}{rgb}{0.858, 0.188, 0.478}
%\definecolor{olive}{RGB}{85, 107, 47}
\definecolor{orange}{RGB}{204, 85, 0}

%\definecolor{mypink3}{cmyk}{0, 0.7808, 0.4429, 0.1412}
%\definecolor{mygray}{gray}{0.6}
%\definecolor{olivegreen}[RGB]{85, 107, 47}

\newcommand{\K}{K}
\newcommand{\M}{M}
\newcommand{\Mo}{M_o}
\newcommand{\So}{S_o}
\newcommand{\Smod}{S}
\newcommand{\parcolor}[1]{{\color{orange}#1}}
\newcommand{\Ts}{T_s^{\rm MPC}}
\newcommand{\cites}[1]{\begin{small}(#1)\end{small}}
%\newcommand{\Ts}{T_s}

\definecolor{darkgreen}{RGB}{20,150,50}
\definecolor{acq}{RGB}{119, 172, 48}
\definecolor{true}{RGB}{162, 20, 47}
\definecolor{surr}{RGB}{0, 114, 189}

\begin{document}

\begin{frame}
  \titlepage
\end{frame}

\begin{frame}{Motivations}
Calibration of \structure{model predictive controllers} is \alert{costly} and \alert{time-consuming}.\\
One has to choose model, cost function weights, prediction/control horizon, design a state estimator, etc.
%Typically, models are obtained through Physical modeling or Identification
\vskip 1em
The design is even more involved for \structure{embedded applications}, where the hardware is also a limit. Need to choose carefully:
\begin{itemize}
\item Solution strategy (QP solver?) 
\item Low-level solver settings (e.g., tolerances)
\item A \structure{feasible} MPC sampling time
\end{itemize}
%\vskip 1em
%Requires \structure{knowledge} in control systems, numerical math, computer engineering + \structure{hands-on experience}.
\vskip 1em
In this paper, we introduce an approach to ease embedded MPC calibration based on \structure{global optimization} and repeated experiements.
\end{frame}


\begin{frame}{Control architecture}{Model Predictive Controller}
\begin{figure}
	\centering
		\includegraphics[width=.3\textwidth]{img/MPCloop.pdf}
\end{figure}

Linear constrained MPC for a (non-linear) plant with model
\begin{columns}
\column{.3\textwidth}
\begin{small}
\begin{equation*}
\begin{split}
\dot x(t) &= f(x(t), u(t);\; \parcolor{\beta})\\
y(t)  &=g(x(t); \; \parcolor{\beta}) %+\parcolor{D}u_t,
\end{split}
\end{equation*}
\end{small}
\column{.2\textwidth}
$$\Rightarrow$$
\column{.5\textwidth}
\begin{small}
\begin{equation*}
\begin{split}
x_{t+1}&={A}(\parcolor{\Ts}, \parcolor{\beta})x_t+{B}(\parcolor{\Ts}, \parcolor{\beta})u_t\\
y_t &=C(\parcolor{\beta})x_t %+\parcolor{D}u_t,
\end{split}
\end{equation*}
\end{small}
\end{columns}
input $u$ chosen on-line as the solution (applied in 
{receding horizon}) of:
\vskip -1.5em
\begin{small}
\begin{align*}
&\min_{\left\{u_{t+k|t}\right\}_{k=1}^{\parcolor{N_c}}} \;\;\sum_{k=1}^{\parcolor{N_p}}  \norm{\hat y_{t+k|t}-\overline{y}_{t+k}}^2_{\parcolor{Q_y}}
 + \norm{u_{t+k|k}\!-\!u_{t+k-1|t}}^2_{\parcolor{Q_{\Delta u}}}  \\ 
&  \text{s.t. model equations + constraints on  $y$, $u$, $\Delta u$}    
\end{align*}
\end{small}
%\vskip -1.5em
%MPC control law computed on-line using a QP solver.% (OSQP, ADMM-based).
\end{frame}

\begin{frame}{Control architecture}{Design variables}
The MPC design variables to be tuned are
\begin{columns}[t]
\column{.6\textwidth}
\begin{itemize}
\item Model parameters: $\parcolor{\beta}$
\item MPC sampling time: $\parcolor{\Ts}$
\item MPC cost function: $\parcolor{Q_y}, \parcolor{Q_{\Delta u}}, \parcolor{N_c},  \parcolor{N_p}$
\end{itemize}
\column{.4\textwidth}
\vskip -1em
\begin{figure}
	\centering
		\includegraphics[width=.75\textwidth]{img/MPCloop.pdf}
\end{figure}
\end{columns}
\begin{itemize}
\item MPC implementation: QP solver tolerances $\parcolor{\epsilon_{\rm abs}}, \parcolor{\epsilon_{\rm rel}}$
\item State estimator: Kalman covariances $\parcolor{W_v}, \parcolor{W_w}$
\end{itemize}
%\vskip 1em
%in order to: 
%\begin{enumerate}
%\item maximize closed-loop performance %objective $J$
%\item ensure that the implemented MPC runs in real-time
%\end{enumerate}
%\vskip 1em
%Note: closed-loop performance is not the MPC cost function!
\end{frame}



\begin{frame}{MPC calibration procedure}{Overview}
In our optimization-based  MPC calibration, we  define
\begin{itemize}
\item Tunable \structure{design parameters} collected in a \parcolor{design vector $\theta$} $\in \Theta$.
\item A closed-loop \structure{performance index $J$} defined in terms of measured input/outputs during the {calibration experiment}: $J = J(y_{1:T},u_{1:T};\;\parcolor{\theta})$
\item An \structure{procedure} to perform \structure{calibration experiments} (or SIL simulations) representative of the intended closed-loop operation
\end{itemize}
\pause
MPC calibration is seen as a \structure{global optimization} problem:
\begin{align*}
&\theta^{\rm opt} = \argmin_{\theta \in \Theta} J(y_{1:T},u_{1:T};\;{\theta})\\
&  \text{s.t. } T_{\rm calc}^{\rm MPC}({\theta}) \leq {\Ts} 
\end{align*}
each (noisy) \structure{function evaluation}  correspond to a \structure{calibration experiment}.\\
Both $J$ and $T_{\rm calc}^{\rm MPC}$ (worst-case MPC computation time) are observed!
%\vskip 1em
%\pause
%Problem is tackled using \structure{efficient} global optimization algorithms.
\end{frame}

\begin{frame}{Global Optimization Algorithm}{Overview}
 Several \structure{derivative-free} global optimization algorithms exist:
 \begin{itemize}
  \item Response surface methods
  \item Bayesian Optimization% (BO)
  \item Genetic algorithms
  \item Particle Swarm Optimization
  \item \dots%Lipshitzian-based partitioning techniques
  \end{itemize}  
% \end{itemize}
 \vskip .5em
 \pause
 Here, we adopt \structure{GLIS}: a method recently introduced in \cites{Bemporad, 2019}.
 \begin{itemize}
  \item Radial basis function \structure{surrogate} + inverse distance weighting (IDW) for \structure{exploration}
  \item Performs better than BO on standard benchmarks, at a lower computational cost
  %\item Matlab- and Python-based implementation of GLIS available
 \end{itemize}
\end{frame}

\begin{frame}{GLIS Algorithm}{Surrogate function}

\begin{columns}[t]
 \column{.55\textwidth}
  Goal: solve \structure{global optimization} problem:
% \begin{align*}
$$  \min_{\theta \in \Theta} J(\theta)$$
%\\
%&  \text{s.t. } \theta \in \Theta 
%\end{align*}
%\vskip -.5em
\begin{enumerate}
\item Collect ${n_{\rm in}}$ initial observations $\{(\theta_1, J_1), (\theta_2, J_2), \dots, (\theta_{n_{\rm in}}, J_{n_{\rm in}})\}$
\item Build a \structure{surrogate function}
\vskip -0.5em
$$\hat J(\theta) = \sum_{i=1}^{n} \alpha_i \phi(\norm{\theta - \theta_i}_2 )$$
\end{enumerate}
 \column{.45\textwidth}
 %\vskip -1.5em
 	\centering
 		True {\color{true} $J(\theta)$}, surrogate {\color{surr} $\hat{J}(\theta)$}
 \begin{figure}
		\includegraphics[width=.7\textwidth]{img/GLIS/surrogate.png}
\end{figure}
$ \phi = $ \text{radial basis function}\\
\vskip .5em
Example: $\phi(d) = \frac{1}{1 + (\epsilon d)^2}$
\end{columns}
\vskip 1em
%Vector $\alpha$ solves $\hat J(\theta_i) = J_i,\; i=1,\dots,n\;\;$ (= linear system) 
\pause
\vskip 1em
Minimizing $\hat J(\theta)$ to find $\theta_{n+1}$ may easily miss the global optimum! 
\end{frame}

\begin{frame}{GLIS Algorithm}{Exploration vs. Exploitation}
\begin{columns}[t]
 \column{.57\textwidth}
\begin{enumerate}
\setcounter{enumi}{2}
\item Construct the IDW \structure{exploration function} 
\begin{equation*}
 z(\theta) = \frac{2}{\pi}\Delta F \tan^{-1}\left(\frac{1}{\sum_{i=1}^n w_i(\theta)} \right)
\end{equation*}
where $w_i(\theta) = \frac{e^{-\norm{\theta - \theta_i}^2}}{\norm{\theta - \theta_i}^2}$%{-\norm{x - x_i}}^2%}{1}$
\item Optimize the \structure{acquisition function}:
$$\theta_{n+1} = \arg \min_{\theta \in \Theta} \hat J(\theta) - \alert{\delta} z(\theta)$$
\end{enumerate}
 \column{.45\textwidth}
 	\centering
 \begin{figure}
	Exploration function {\color{acq}$z(\theta)$}
		\includegraphics[width=.7\textwidth]{img/GLIS/acquisition.png}
\end{figure}
$\alert{\delta}$: exploitation vs.\\ exploration trade-off
\end{columns}
\vskip 1em
to get the \structure{query point} $\theta_{n+1}$. Execute experiment with $\theta_{n+1}$, measure $J_{n+1}$.
\pause
\vskip 1em
Iterate the procedure for $n+2, n+3 \dots$.  
\end{frame}

\begin{frame}{MPC Calibration Example}%{Cart-pole system}
\begin{columns}
\column{.48\textwidth}
\centering
Cart-pole system
\begin{figure}
    \includegraphics[width=.45\textwidth]{img/example/cart_pole.pdf}
\end{figure}
    \begin{itemize}
     \item Input: force $F$ with \structure{fast disturbance} ($5~\text{rad/sec}$)
     \item Output: noisy  position $p$%$ with white measurement noise
     \item Objective: follow trajectory for $p$, keep $\phi$ small.
     \end{itemize}
\column{.54\textwidth}
Optimization-based calibration of
\begin{enumerate}
 \item MPC sample time $T_{s}^{\rm MPC}$
 \item MPC weights $Q_y$ and $Q_{\Delta u}$
 \item Prediction and control horizon
 \item Kalman filter covariance matrices
 \item QP solver's abs. and rel. tolerances
 \end{enumerate}

\vskip 1em
 $n_\theta=14$ design parameters optimized using $n=500$ iterations of GLIS

 \begin{small}
 \begin{equation*}
 J =  \int_0^{T_{\rm exp}} |\overline{p}(t) - p(t)| +  30 |\phi(t)|\;dt 
\end{equation*}
\end{small}
% \footnotesize
% \url{https://github.com/forgi86/efficient-calibration-embedded-MPC.git}
\end{columns}
%\pause
%\vskip 1em
\pause
 Python codes on-line, $100\%$ \structure{open source} (including all dependencies).\\
\small
\url{https://github.com/forgi86/efficient-calibration-embedded-MPC}
\end{frame}

\begin{frame}{MPC Calibration Example}%{Cart-pole system}
\begin{columns}
\column{.3\textwidth}
\begin{figure}
	\centering
		%GLIS iterations on the PC
		\includegraphics[width=.99\textwidth]{img/example/HIL_PC.pdf}
\end{figure}
\column{.2\textwidth}
\centering
\huge vs.
\column{.3\textwidth}
\begin{figure}
		%GLIS iterations on the Raspberry PI
	\centering
    \includegraphics[width=.99\textwidth]{img/example/HIL_PI.pdf}
\end{figure}
\end{columns}
\vskip 1em
\begin{itemize}
 \item An Intel i5 PC (left) vs. an ARM-based Raspberry PI 3 (right)
 %\item On the same control problem, running the same Python code
 \item PI is about 10 time slower than the PC for MPC computations
 %\item Cart-pole simulated in the same Python code for simplicity \dots
 %\item Could be used in a true HIL setup or a real system
\end{itemize}
\end{frame}


\begin{frame}{MPC Calibration Example}%{Cart-pole system}
\begin{columns}
\column{.5\textwidth}
\begin{figure}
	\centering
		Optimized MPC on the PC
		\includegraphics[width=.99\textwidth]{img/example/BEST_GLIS_PC.pdf}
		$T_s^{\rm MPC} = 6$~ms
\end{figure}
\column{.5\textwidth}
\begin{figure}
	Optimized MPC on the Raspberry PI 
	\centering
    \includegraphics[width=.99\textwidth]{img/example/BEST_GLIS_PI.pdf}
    $T_s^{\rm MPC} = 22$~ms
\end{figure}
\end{columns}
\vskip 1em
\begin{itemize}
 \item Position and angle control tighter on the PC
 \item Faster loop update on the PC $\Rightarrow$ more effective disturbance rejection
 \item Calibration squeezes max performance out of the hardware!
\end{itemize}
\end{frame}

\begin{frame}{MPC Calibration Example}%{Cart-pole system}
\begin{columns}
\column{.5\textwidth}
\begin{figure}
	\centering
		GLIS iterations on the PC
		\includegraphics[width=.99\textwidth]{img/example/ITER_GLIS_PC_500.pdf}
\end{figure}
\column{.5\textwidth}
\begin{figure}
		GLIS iterations on the Raspberry PI
	\centering
    \includegraphics[width=.99\textwidth]{img/example/ITER_GLIS_PI_500.pdf}
\end{figure}
\end{columns}
\vskip 1em
\begin{itemize}
 \item High cost for interrupted tests/unfeasible design
 \item For increasing $n$, more points have ``low'' cost
 \item Experiments with high cost still appear for large $n$ (exploration)
 %\item Optimum on PC is slightly better than on the PI. Some configurations are only feasible on PC!
\end{itemize}
\end{frame}

\begin{frame}{Simulation Example}{Cart-pole system}
Similar calibration results are obtained using Bayesian Optimization (BO).\\
However, BO is a much heavier algorithm!
\vskip -.5em
 \begin{columns}
\column{.5\textwidth}
\begin{figure}
	\centering
		GLIS iterations on the PI
		\includegraphics[width=.99\textwidth]{img/example/COMPUTATION_GLIS_PI.pdf}
\end{figure}
\column{.5\textwidth}
\begin{figure}
	BO iterations on the PI
	\centering
		\includegraphics[width=1\textwidth]{img/example/COMPUTATION_BO_PI.pdf}
\end{figure}
\end{columns}
\begin{itemize}
\item With GLIS, the overall computation time is dominated by function evaluations (= running MPC and simulating the system).
%\item A lightweight algorithm like GLIS could be used for self-tuning on an embedded system!
\end{itemize}
\end{frame}

\begin{frame}{Conclusions}
An approach for embedded MPC calibration approach based on global optimization and repeated experiments
\begin{itemize}
\item Higher- and lower-level tuning knobs jointly optimized for performance, while keeping the design feasible for real-time implementation.
\item Tested in simulation on two hardware platforms
\end{itemize}
\vskip 1em
\pause
Current/future works
\begin{itemize}
\item Application to {robotic systems}
\item Preference-based MPC calibration
\end{itemize}
\end{frame}

\begin{frame}{}{}
\begin{center}
\huge{\structure{Thank you.\\ Questions?}}\\
\vskip 1em
\begin{small}
\texttt{marco.forgione@idsia.ch}
\end{small}
\end{center}
\end{frame}

%\appendix
%\newcounter{finalframe}
%\setcounter{finalframe}{\value{framenumber}}
%\setcounter{framenumber}{\value{finalframe}}
	
\end{document}
